\begin{papercard}{hallgren2012}
\cf{Stages} \stagetag{I}. A tutorial on computing and \emph{reporting} inter-rater reliability (IRR): design choices, common mistakes, Cohen's $\kappa$ and its variants for nominal data, intraclass correlations (ICCs) for ordinal to ratio data, with SPSS and R syntax and two model reporting paragraphs. The most practical, balanced citation for choosing among $\kappa$, $\alpha$ and ICC.

\cf{Concept} IRR quantifies agreement between two or more coders making independent ratings. In classical test theory $X=T+E$ and IRR is the share of observed variance due to true-score variance (an IRR of 0.80 means 80\,\% true variance) \pc{hallgren2012}{2}{IRR analysis aims to determine how much of the variance in the observed scores is due to variance in the true scores}. IRR is distinct from validity: an instrument can have good IRR and poor validity.

\cf{Design decisions to make a priori} (1) Do all subjects get multiple coders, or a subset (the IRR of the subset generalises)? (2) \emph{Fully crossed} (the same coders rate all) or not; crossed designs allow coder bias to be modelled \pc{hallgren2012}{2}{Although fully crossed designs can require a higher overall number of ratings to be made, they allow for systematic bias between coders to be assessed}. (3) \emph{Psychometric properties}: restriction of range lowers IRR even for a good scale, so consider changing the scale and pilot it \pc{hallgren2012}{2}{Restriction of range often lowers IRR estimates because the component of equation 3 is reduced}. (4) Coder training: set and report an a priori IRR level to reach before coding the real subjects \pc{hallgren2012}{3}{In these cases it is common to specify an a priori level of IRR that must be achieved before subjects from the real study are rated}.

\cf{Common mistakes} (1) Percent agreement ignores chance (two coders guessing 50/50 agree 50\,\%; guessing 10/90 agree 82\,\% by chance alone) \pc{hallgren2012}{3}{if coders were to randomly rate 50\% of subjects as “depressed” and 50\% as “not depressed” without regard to the subject’s actual characteristics}. (2) Not reporting the statistic and its \emph{variant}. (3) Wrong statistic for the design or metric. (4) Not computing IRR on variables in their \emph{final transformed form}. (5) Not interpreting the effect of IRR on statistical power: low IRR adds noise and raises type-II error \pc{hallgren2012}{4}{Low IRR may increase the probability of type-II errors}. Report possible reasons (restricted range, poor scale, poorly trained coders, hard-to-observe construct) and decisions about dropping variables \pc{hallgren2012}{4}{Possible reasons for low IRR should be discussed}.

\cf{Kappa for nominal variables} $\kappa=(P(a)-P(e))/(1-P(e))$ with chance agreement from the coders' marginal frequencies. Worked example: $P(a)=.79$, $P(e)=.50$, $\kappa=.58$ \pc{hallgren2012}{5}{The total probability of any chance agreement would then be 0.225 + 0.275 = 0.50}. Interpretation: Landis \& Koch (0.61--0.80 substantial), versus Krippendorff's more conservative cut-offs (below .67 discount; .67--.80 tentative; above .80 definite), which are often not applied in practice \pc{hallgren2012}{5}{Krippendorff (1980) provides a more conservative interpretation suggesting that conclusions should be discounted for variables with values less than 0.67}. \textbf{Prevalence problem}: very unequal category frequencies make $\kappa$ unrepresentatively \emph{low}; \textbf{bias problem}: different marginals between coders make it unrepresentatively \emph{high}. Variants: Byrt et al.\ (prevalence-adjusted), Siegel \& Castellan (bias-adjusted, the only one in SPSS); no single variant corrects both, so report several if needed \pc{hallgren2012}{5}{No single kappa variant corrects for both bias and prevalence}. Weighted $\kappa$ (Cohen 1968) for ordered categories; quadratic-weighted $\kappa$ equals a two-way mixed consistency ICC. For three or more coders: Fleiss' $\kappa$ (a random sample of coders per subject, so inappropriate for fully crossed designs), or Light's average of pairwise $\kappa$, or Davies \& Fleiss \pc{hallgren2012}{6}{For fully-crossed designs with three or more coders, Light (1971) suggests computing kappa for all coder pairs}.

\cf{Worked kappa example and report template} Three coders, 12 subjects, three categories. Check marginals for prevalence/bias, compute pairwise $\kappa$ (0.62, 0.61, 0.80), average 0.68; SPSS's Siegel--Castellan version averages 0.56 \pc{hallgren2012}{6}{The resulting estimate of Cohen’s kappa averaged across coder pairs is 0.68}. Report: the variant chosen and why (marginals did not show prevalence or bias problems), the averaging method, the qualitative interpretation, comparison with prior studies, and the implication for power \pc{hallgren2012}{7}{Reporting of these results should detail the specifics of the kappa variant that was chosen}.

\cf{ICCs} For ordinal, interval and ratio data with two or more coders; they take the \emph{magnitude} of disagreement into account. Choose along four dimensions: one-way vs.\ two-way model (fully crossed $\Rightarrow$ two-way); \emph{absolute agreement} vs.\ \emph{consistency}; \emph{single-} vs.\ \emph{average-measures} (average values are higher); random vs.\ mixed effects (same numbers, different generalisation) \pc{hallgren2012}{8}{Four major factors determine which ICC variant is appropriate based on one’s study design}. ICCs use list-wise deletion, so with many missing values Krippendorff's $\alpha$ may fit better \pc{hallgren2012}{8}{ICCs use list-wise deletion for missing data}. Cicchetti's bands: $<.40$ poor, .40--.59 fair, .60--.74 good, $\ge.75$ excellent. Example: ten subjects, three coders, two-way mixed consistency average-measures ICC $=.96$ \pc{hallgren2012}{9}{The resulting ICC is high, ICC = 0.96, indicating excellent IRR for empathy ratings}. Significance tests of ICC $=0$ are rarely reported.

\cf{Conclusion} Use validated IRR statistics, never percent agreement alone, and fit the statistic to design and distributions. $\alpha$ is more flexible with missing data and across metrics but less well known \pc{hallgren2012}{10}{Krippendorff’s alpha can be generalized across nominal, ordinal, interval, and ratio variable types}.

\cf{Use for this thesis} A reporting checklist: design (fully crossed? subset?), variable metric per code, the exact statistic and variant, marginal distributions (prevalence/bias), per-pair and averaged values, interpretation scale chosen \emph{before} analysis, and effects on later analyses. For a binary code set coded by two coders on the same sample (fully crossed), report $\kappa$ \emph{and} prevalence-adjusted variants \emph{and} $\alpha$ with their marginals, and say why they differ if they do. It is a tutorial, not a decision on which statistic is best, so pair it with Gwet, Feinstein \& Cicchetti and Krippendorff for the paradox debate, and with the thesis' own pre-stated threshold.
\end{papercard}
